\section*{Chapitre \ref{ch:ll:measuring-latency} --- Mesurer la latence}

\begin{solution}{exo:ll:measuring-latency:1}
$(3\,020 - 45)/3.02 = \qty{985.1}{\nano\second}$.
\end{solution}

\begin{solution}{exo:ll:measuring-latency:2}
1\,000 a une longueur de bits égale à 10, donc le décalage vaut 2 et les bits de tête valent $\lfloor 1\,000/4 \rfloor = 250$~: l'indice $256 + 128 +
(250-128) = 506$, couvrant de 1\,000 à 1\,003.
\end{solution}

\begin{solution}{exo:ll:measuring-latency:3}
$1\,024 + 53 \times 512 + 512 = 28\,672$ compteurs, avec une précision relative de $2^{-9} \approx 0.195~\%$.
\end{solution}

\begin{solution}{exo:ll:measuring-latency:4}
7, 5 et \qty{3}{\milli\second}~: les valeurs $9 - 2k$ tant qu'elles dépassent l'intervalle de \qty{2}{\milli\second}.
\end{solution}

\begin{solution}{exo:ll:measuring-latency:5}
La correction suppose que chaque requête manquante aurait reçu une réponse dès la fin du gel. En réalité, les dix
requêtes programmées pendant le gel se mettent en file les unes derrière les autres et sont servies à la suite, si bien que celles envoyées plus tard
pendant le gel attendent aussi les précédentes~; le testeur en boucle ouverte mesure cette file, la correction non.
\end{solution}

\begin{solution}{exo:ll:measuring-latency:6}
Il n'enregistre que leur temps de traitement, comme s'ils étaient arrivés au moment où la boucle les atteint~: rien de l'attente, qui atteint au plus
\qty{2}{\milli\second} dans la socket. Il faudrait mesurer à partir de l'arrivée du paquet~: un \omterm{def:ll:measuring-latency:hwts}{horodatage matériel} fourni par
la carte, ou au moins l'horodatage de réception du noyau (\texttt{SO\_TIMESTAMPNS}).
\end{solution}

\begin{solution}{exo:ll:measuring-latency:7}
\texttt{ll\_measure.spectrum(stall=50\_000, period=5e6)}~: boucle fermée \qty{0.16}{\milli\second}, corrigée
\qty{1.07}{\milli\second}, boucle ouverte \qty{5.77}{\milli\second}. Le testeur en boucle fermée prélève un échantillon par gel, désormais
20 sur 99\,000 (0.02~\%), très en dessous de 1~\%~: son p99 ne voit pas du tout les gels, alors que les utilisateurs mis en file derrière chaque
gel de \qty{50}{\milli\second} les voient.
\end{solution}

\begin{solution}{exo:ll:measuring-latency:8}
Trois défauts~: un seul client en boucle fermée (l'\omterm{def:ll:measuring-latency:co}{omission coordonnée} masque chaque arrêt)~; des seaux linéaires de \qty{100}{\nano\second}
jusqu'à \qty{100}{\micro\second} ne peuvent pas du tout contenir la queue, si bien que toute valeur au-delà du dernier seau est perdue ou ramenée à la borne~; et
un million d'échantillons ne donnent que 100 valeurs au-delà du p99.99, trop peu pour l'affirmer sans une dispersion. Il faut utiliser un générateur en boucle ouverte au
débit de production, un histogramme log-linéaire et des séries répétées.
\end{solution}

\begin{solution}{pb:ll:measuring-latency:1}
\begin{enumerate}
\item Le testeur en boucle ouverte envoie 100\,000 requêtes~; celui en boucle fermée en envoie 99\,000, car chaque gel avale les dix
millisecondes pendant lesquelles il attend une réponse.
\item 100 sur 99\,000~: 0.10~\%.
\item Les dix requêtes programmées pendant chaque gel, soit 1~\% du total, plus celles mises en file derrière elles juste après.
\item Neuf par gel (\qty{9.1}{\milli\second}, \qty{8.1}{\milli\second}, \dots, \qty{1.1}{\milli\second})~: 900, de sorte que l'histogramme corrigé contient
99\,900 échantillons.
\item \qty{0.16}{\milli\second} (boucle fermée), \qty{1.06}{\milli\second} (corrigé) et \qty{1.87}{\milli\second} (boucle
ouverte).
\item À la médiane, tous les testeurs voient la réponse ordinaire de \qty{100}{\micro\second}~; au-delà du p99.9, les trois voient le gel
lui-même, soit environ \qty{10.1}{\milli\second}.
\item La correction suppose que les requêtes manquantes auraient été servies à la fin du gel~; la file qui
se forme entre-temps les fait attendre plus longtemps.
\item Le chiffre en boucle ouverte~: les utilisateurs envoient selon leur propre calendrier.
\item $2^{-7} \approx 0.8~\%$~; sans importance ici, où les testeurs diffèrent d'un facteur dix.
\item 7\,424.
\item Mécaniquement oui, en additionnant les effectifs~; mais fusionner les sorties de testeurs différents mélange des questions différentes. Il faut
fusionner les histogrammes d'un même testeur entre les séries ou les machines.
\item Chaque échantillon de gel, vers \qty{10.1}{\milli\second}, tomberait au-delà du dernier seau et serait ramené à la borne ou perdu~:
le rapport n'aurait pas de queue.
\item \textbf{Résultat nommé.} Le testeur en boucle fermée donne un p99 de \qty{0.16}{\milli\second} brut et
\qty{1.06}{\milli\second} corrigé~; le testeur en boucle ouverte, qui mesure ce que vivent les utilisateurs, donne \qty{1.87}{\milli\second}.
\item Instants~: de l'heure d'envoi prévue à la réception de la réponse côté client. Charge~: boucle ouverte, 1\,000 requêtes par
seconde pendant 100 secondes. Histogramme~: log-linéaire, précision de 0.8~\%, 100\,000 échantillons~; p50 \qty{0.10}{\milli\second}, p99
\qty{1.87}{\milli\second}, p99.9 \qty{10.1}{\milli\second}. Il faut ajouter la machine, le placement, les versions, l'horloge et les répétitions.
\item La collecte complète périodique d'un ramasse-miettes (garbage collector), une rotation de journaux, une tâche cron, une tempête d'interruptions de minuterie~: tout travail périodique
qui interrompt le service pendant quelques millisecondes.
\item Chercher un générateur en boucle fermée (une seule requête en attente par client), un nombre d'échantillons inférieur au débit prévu multiplié
par la durée, ou un p99.9 très supérieur à un p99 plat.
\item Une stratégie qui horodate les messages au moment où sa boucle les atteint, plutôt qu'au moment où ils sont arrivés.
\item Parce qu'elle est présente dans chaque figure~: une horloge de \qty{20}{\nano\second} double une étape de \qty{20}{\nano\second}.
\item Le temps écoulé entre l'arrivée du paquet sur le câble et la première lecture de celui-ci par le programme~: le chemin de réception et toute
mise en file dans la socket, que les horodatages logiciels ne voient pas.
\item Les deux instants, qui les a horodatés, la charge, les percentiles avec leur effectif, et les conditions.
\end{enumerate}
\end{solution}

\begin{solution}{iq:ll:measuring-latency:1}
On lit le \omterm{def:ll:cpu-microarchitecture:tsc}{compteur d'horodatage} (\texttt{rdtsc}, ou \texttt{rdtscp} pour attendre les instructions précédentes) et l'on convertit avec un taux
calibré par rapport à \texttt{CLOCK\_MONOTONIC\_RAW}. Il n'a de sens d'un cœur à l'autre que si le compteur est invariant
(cadence constante dans tous les états de puissance, indicateurs \texttt{constant\_tsc} et \texttt{nonstop\_tsc}) et synchronisé entre
les sockets, ce que le noyau vérifie au démarrage.
\iqlookfor{TSC, invariance, calibration par rapport à une horloge brute.}
\end{solution}

\begin{solution}{iq:ll:measuring-latency:2}
Un générateur qui attend chaque réponse cesse d'envoyer pendant les arrêts, si bien que les requêtes qui auraient attendu ne sont jamais
mesurées. Il faut piloter la charge en boucle ouverte, selon un calendrier indépendant des réponses, et mesurer à partir de l'heure d'envoi prévue de
chaque requête~; ou bien enregistrer avec l'intervalle attendu pour corriger.
\iqlookfor{le mécanisme de la boucle fermée et le remède de la boucle ouverte.}
\end{solution}

\begin{solution}{iq:ll:measuring-latency:3}
Un histogramme log-linéaire à seaux fixes (exact en dessous de $2^S$, puis $2^{S-1}$ seaux par puissance de deux)~: mémoire constante,
quelques instructions par enregistrement, erreur relative bornée, fusionnable entre les threads~; on lit le percentile par le rang le plus proche
sur les effectifs.
\iqlookfor{erreur relative bornée et mémoire fixe, et non des échantillons stockés.}
\end{solution}

\begin{solution}{iq:ll:measuring-latency:4}
La moyenne mélange le corps et la queue et les masque l'un et l'autre~; quelques arrêts la déplacent peu alors qu'ils coûtent le plus. Il faut indiquer le p50,
le p99, le p99.9, le maximum et l'effectif, à partir d'un histogramme, avec les conditions.
\iqlookfor{percentiles et effectif.}
\end{solution}

\begin{solution}{iq:ll:measuring-latency:5}
Utile pour repérer où part le temps moyen dans un cas lent reproductible (une fonction chaude, des \omterm{def:ll:memory-hierarchy-and-caches:miss}{défauts de cache} (cache miss) avec les
compteurs matériels)~; inutile pour les événements rares de queue, que le profileur échantillonne rarement, et pour tout ce que l'interruption d'échantillonnage perturbe.
\iqlookfor{moyenne contre queue.}
\end{solution}

\begin{solution}{iq:ll:measuring-latency:6}
Plusieurs séries alternées de chaque version sur la même machine calme et épinglée~; la dispersion du p99 d'une série à l'autre~; une
comparaison statistique (bootstrap ou permutation des séries)~; et un mécanisme qui explique l'écart. Le chapitre~25 construit
la barrière.
\iqlookfor{répétition et estimation du bruit avant de conclure.}
\end{solution}
